% Contenido del póster: Vives y Ye (2025), Journal of Financial Economics. DOI 10.1016/j.jfineco.2025.104040
% Elaborado con apoyo de IA a partir del resumen oficial del artículo (OpenAlex), 2026-11-12.
\def\DOI{10.1016/j.jfineco.2025.104040}
\def\TituloEs{Entrada de las fintech, competencia en el mercado de préstamos y bienestar}
\def\TituloOriginal{Fintech entry, lending market competition, and welfare}
\def\Autores{Vives y Ye}
\def\Anio{2025}
\def\Revista{Journal of Financial Economics}

\def\Problema{Las fintech entran a los mercados de préstamos aunque supervisan peor a los prestatarios que los bancos. Hace falta explicar por qué lo logran y qué consecuencias tiene para la inversión y el bienestar.}
\def\Objetivo{Analizar la competencia entre bancos y fintech en el mercado de préstamos y su efecto en la inversión y el bienestar.}
\def\Metodo{Marco espacial de competencia construido sobre las diferencias clave entre bancos y fintech, cuyos resultados se contrastan con la evidencia empírica disponible.}
\def\Resultados{\begin{itemize}
  \item Las fintech pueden entrar aunque supervisen peor, gracias a su mayor flexibilidad para fijar precios.
  \item Una mayor concentración bancaria se asocia con un mayor volumen de préstamos fintech.
  \item Con costos de fondeo similares, los prestatarios fintech pagan menos y tienen más impagos; si el fondeo fintech es mucho más caro, ocurre lo contrario.
  \item La conveniencia que ofrecen las fintech también les permite cobrar más que los bancos.
  \item La entrada mejora el bienestar si las fintech supervisan bien y la competencia entre ellas es intermedia.
\end{itemize}}
\def\Conclusion{El efecto de las fintech sobre tasas, inversión y bienestar no es único: depende de su costo de fondeo, de su capacidad de supervisión y de cuánto compiten entre sí.}
\def\Relevancia{El estudio muestra por qué el precio del crédito de plataforma puede ser mayor o menor que el bancario. En el Perú no puede comprobarse cuál es el caso, porque la tasa de cada operación no se publica. Además, la oferta peruana está muy concentrada: el índice de Herfindahl-Hirschman del financiamiento participativo llegó a \cifra{hhi2025} en 2025.}
\def\PalabrasClave{fintech, competencia, concentración bancaria, costo de fondeo, bienestar}
% --- Ficha del estudio (herramientas/posteres/ficha_articulos.py) ---
\def\Ambito{Teórico}
\def\Enfoque{Modelo espacial de competencia entre bancos y fintech}
\def\FiguraPeru{1}
% --- fin de la ficha ---
